\label{sec:justificacion}

Las redes neuronales informadas por física (PINN, por sus siglas en inglés) arrastran una limitación que no se resuelve agregando parámetros \cite{grossmann2024,krishnapriyan2021}. En las arquitecturas basadas en perceptrones multicapa (MLP, por sus siglas en inglés), el entrenamiento privilegia las componentes de baja frecuencia y aprende con dificultad las de alta frecuencia o multiescala \cite{rahaman2019,tancik2020,wang2021eigenvector}. Esta tendencia es especialmente relevante en ecuaciones diferenciales parciales (PDE, por sus siglas en inglés), donde una solución puede presentar bajo error agregado y aun así representar de forma deficiente regiones localizadas de alta variación física \cite{krishnapriyan2021,deryckmishra2022,grossmann2024}. El problema central, entonces, está menos en el tamaño del modelo que en si la arquitectura puede representar con fidelidad el contenido local que gobierna la solución.

El peso de esta limitación crece en los sistemas piezoeléctricos, donde la respuesta mecánica y la eléctrica están acopladas mediante ecuaciones constitutivas y condiciones de frontera \cite{tandis2025piezoelectric,mouratidou2026shunted}. En estos problemas, la solución puede concentrar información físicamente relevante en interfaces, discontinuidades geométricas, zonas de concentración de esfuerzo o regiones de alto gradiente \cite{wang2004intensity}. Son justamente esas regiones las que condicionan la validez de la respuesta del sistema. Cuando la representación aprendida pierde fidelidad en esos puntos, el modelo no solo degrada su precisión local: también compromete la interpretación física de la solución.

El problema es también de diagnóstico. Buena parte de las evaluaciones de PINN reporta error relativo global, residual promedio o comparaciones agregadas contra soluciones numéricas \cite{raissi2019,lu2021deepxde,grossmann2024}. Estas métricas son útiles para comparar modelos, pero reducen la distribución espacial del error a un solo valor. Un modelo puede entonces mejorar su desempeño promedio sin tocar las zonas donde el error pesa más en términos físicos \cite{krishnapriyan2021}. En ese escenario, una mejora métrica no garantiza una mejora física.

En los sistemas acoplados esto tiene una consecuencia adicional. Cuando una variable aparece en las ecuaciones que gobiernan las otras, el error local puede afectar la consistencia del sistema completo. Una falla en el campo mecánico puede alterar la estimación del campo eléctrico, y una falla en el campo eléctrico puede modificar la respuesta mecánica \cite{mouratidou2026shunted,tandis2025piezoelectric}. La representación en regiones de alto gradiente deja entonces de ser un problema puntual y pasa a ser una condición para mantener la coherencia entre campos.

Las estrategias adaptativas actuales atacan parte de esta dificultad desde la distribución del entrenamiento. El refinamiento adaptativo basado en residuales (RAR, por sus siglas en inglés) agrega puntos donde el residual es alto, y las redes neuronales variacionales informadas por física con refinamiento hp (hp-VPINN, por sus siglas en inglés) introducen descomposición de dominio y espacios de prueba locales \cite{lu2021deepxde,wu2023sampling,kharazmi2021hpvpinns}. Estas estrategias son relevantes porque mejoran el muestreo o la formulación variacional, pero no eliminan por sí solas la interrogante sobre si la arquitectura tiene la capacidad interna de representar contenido local de alta frecuencia, especialmente en regiones donde la complejidad física se concentra \cite{rahaman2019,tancik2020,shukla2024fair}. El espacio que persiste no es únicamente computacional: es de representación.

Este trabajo se sitúa en este espacio. La pregunta de fondo no es si una arquitectura alternativa reduce el error global, sino si ofrece un mecanismo interno para asignar capacidad de representación con sensibilidad local. Las redes Kolmogorov-Arnold (KAN, por sus siglas en inglés) y sus variantes informadas por física han abierto una línea reciente para estudiar representaciones alternativas a los MLP en problemas gobernados por ecuaciones diferenciales \cite{liu2024kan,toscano2025pikans,shukla2024fair,wang2025kinn}. En el caso piezoeléctrico, responderla implica evaluar si la solución aprendida conserva validez física en las regiones donde el acoplamiento, la variación multiescala y los gradientes elevados vuelven más probable el fallo de los modelos basados en MLP. La contribución esperada, por tanto, no es una comparación de desempeño, sino una explicación de cómo la representación arquitectónica afecta la consistencia física local de las soluciones aprendidas para PDEs acopladas.